%---------------------------------------------------------------------
\section{The optimal contract and distribution-free decarbonization}\label{sec:optcontract}

Sections~\ref{sec:carbon}--\ref{sec:threshold} signed carbon under the \emph{fixed}
dispatch rule $d^{*}(x)=\min\{(\KF-x)_+,\KM\}$, and found reduction only at the
cheap-renewable corner. The rule, though, was assumed, not optimized. Optimizing it
turns the corner result into a general one and resolves the keystone left open in
\S\ref{sec:threshold}.

\begin{proposition}[The surplus contract is fossil-minimizing; Role~1]\label{prop:surplus}
At any capacities, total fossil in \eqref{eq:profit} is minimized over admissible
dispatch schedules by the \emph{surplus contract} $d^{c}(x)=\min\{(-x)_+,\KM\}$, which
runs the load only on renewable surplus and burns \emph{zero} fossil for the load.
\end{proposition}

Under $d^{c}$ the load's fossil footprint vanishes, so total fossil is the utility's own
output $e(\KR)\eqdef\int_0^{\KF(\KR)}(1-F(x;\KR))\,dx$, which is strictly decreasing in
$\KR$ (more renewables displace fossil). Combined with the capacity channel this gives a
clean, shape-free guarantee.

\begin{theorem}[Distribution-free decarbonization; Role~1]\label{thm:decarb}
Suppose the load is adopted ($\CM<R\,F(0;\KRo)$) and the renewable problem is
single-peaked. Then under the surplus contract $\KRc\ge\KRo$ and
$\E[\text{fossil}^{c}]\le\E[\text{fossil}^{\circ}]$, both strict whenever the load is
adopted---\emph{for every residual-demand law $F$}, with no condition on its shape, tails,
or skew. Equivalently, $\inf_{F\in\Ucal}(\E[\text{fossil}^{\circ}]-\E[\text{fossil}^{c}])\ge0$:
the guarantee is distributionally robust, so the regulator needs no knowledge of $F$.
Relative to no load it is a strict Pareto improvement---profit rises, carbon falls, and
blackout probability is held exactly at the no-load level $1-\varphi^{\circ}$.
\end{theorem}

Weeks~2's finding that the fixed-rule carbon sign is \emph{irreducibly distributional}
(decidable neither way from the cutoffs) is thus a property of the \emph{contract}: the
surplus contract removes the fossil floor that caused it, and the sign becomes definite.

\paragraph{The carbon--profit frontier; Role~1.} The naive and surplus rules are the two
ends of a one-parameter family---grant the load a per-state fossil budget $\tau$,
$d_{\tau}(x)=\min\{(-x)_++\tau,\KM,(\KF-x)_+\}$, with $\tau=0$ the surplus rule and
$\tau\to\infty$ the naive rule. Profit $\Pi(\tau)$ is non-decreasing (proved), and the
exact identity $C(\tau)-C(0)=\Phi(\tau)(1-\Lambda(\tau))$---with $\Phi(\tau)$ the load's
direct footprint and $\Lambda$ the consumer fossil the induced build-out saves per unit
of footprint---reduces global carbon-optimality of the surplus contract to the single
condition $\Lambda\le1$. At interior optima $\Lambda\in[0.14,0.22]$: the induced build-out offsets
roughly a fifth of the load's footprint, not all of it. So the
contracts trace a monotone carbon--profit frontier with no free lunch: every megawatt of
fossil the load is allowed to burn raises emissions. The naive rule is the carbon-worst,
the surplus rule the carbon-best corner; the regulator's reliability and participation
requirements select the point between (the welfare selection is B1's).
